\section{Configuration Format}

The controller reads multi-phase YAML profiles from the \texttt{config/}
folder. Every profile has a \texttt{global} block with the total duration,
warmup and cooldown in seconds, and a list of \texttt{phases}. Each phase has
a name, a duration, an optional \texttt{adaptive\_control} block, and a
\texttt{protocols} block with one entry per protocol that should be active
during that phase. Three profiles ship with the project: \texttt{balanced.yaml}
(constant baseline, then stealth plus random pattern, then a linear ramp on
all four protocols, then adaptive control), \texttt{http2\_heavy.yaml}
(HTTP/2 under load, then an HTTP/2 + QUIC ramp phase, then a burst phase),
and \texttt{mqtt\_heavy.yaml} (a QoS 0 vs.\ QoS 2 vs.\ mixed comparison,
then an MQTT burst phase, then an MQTT ramp phase).

\begin{table}[t]
\centering
\caption{Per-protocol configuration fields}
\label{tab:config-fields}
\footnotesize
\begin{tabular}{@{}p{0.16\linewidth}p{0.78\linewidth}@{}}
\toprule
\textbf{Protocol} & \textbf{Fields} \\
\midrule
http2 & rate, payload\_size, method\_get\_pct, concurrent\_streams, get\_paths, post\_paths, pattern, burst\_rate, burst\_duration, burst\_interval, ramp\_start\_rate, ramp\_end\_rate, ramp\_duration, fault\_rate, extra\_latency\_ms \\
quic & rate, payload\_size, stream\_count, use\_0rtt, pattern, burst\_rate, burst\_duration, burst\_interval, ramp\_start\_rate, ramp\_end\_rate, ramp\_duration, fault\_rate, extra\_latency\_ms \\
mqtt & rate, payload\_size, qos, topic\_count, qos\_distribution, pattern, burst\_rate, burst\_duration, burst\_interval, ramp\_start\_rate, ramp\_end\_rate, ramp\_duration, fault\_rate, extra\_latency\_ms \\
tcp/udp & tcp\_rate, udp\_rate, packet\_size, mode, mean\_interval, min\_size, max\_size, tcp\_ratio, pattern, burst\_size, burst\_interval, ramp\_start\_rate, ramp\_end\_rate, ramp\_duration, fault\_rate, extra\_latency\_ms \\
\bottomrule
\end{tabular}
\end{table}

\texttt{pattern} is the same four-way choice on every protocol: constant,
periodic burst, random (Poisson), or ramp. The first three were already
implemented on at least one generator from the start; periodic burst and
ramp were extended to all four so that any protocol can demonstrate any
sending pattern, not only the ones that happened to need it first. Ramp is
distinct from the global \texttt{warmup}/\texttt{cooldown} period described
above: warmup/cooldown ramps every generator's rate together, once, at the
very edges of a run, whereas \texttt{pattern: ramp} ramps one generator's
rate between \texttt{ramp\_start\_rate} and \texttt{ramp\_end\_rate} for the
duration of whichever phase sets it, independently of every other protocol
active in that phase.

Most field names in the YAML match the generator's Pydantic model exactly
and are forwarded unchanged. Two exceptions exist for readability:
\texttt{method\_distribution: \{GET: 70, POST: 30\}} on the http2 block, and
\texttt{topics: [...]} on the mqtt block. Generators only understand
\texttt{method\_get\_pct} and \texttt{topic\_count} respectively, not a named
distribution or a list of topic names, so the controller's
\texttt{\_phase\_to\_gen\_configs()} translates both before forwarding the
phase config: the GET/POST split is normalized into a single
\texttt{method\_get\_pct} percentage, and the topic list's length becomes
\texttt{topic\_count} (the topic \emph{names} themselves are not
configurable -- the generator only rotates through a fixed built-in set --
so the YAML list is used purely as a readable way to say how many topics
should be active). Any other key that does not match the model is still
silently ignored rather than rejected, since the models accept extra keys
for forward compatibility but the generators only apply the ones they
recognize.

Configuration is not locked to the YAML once the system is running.
\texttt{POST /config/load} switches the active profile and applies phase one
immediately if the system is already started. \texttt{PATCH /generator/\{name\}}
forwards any subset of the fields in Table~\ref{tab:config-fields} to one
generator. The dashboard wraps every one of these fields in a slider, toggle
or text input that calls the same two endpoints, debounced so that dragging a
slider does not flood the controller with requests.
